\chapter{Nichtparametrische Methoden}

\section{Ziel dieses Kapitels}

In den bisherigen Kapiteln haben wir viele parametrische Modelle betrachtet.
Beispiele sind lineare Regression, logistische Regression, neuronale Netze, Gaussian Mixture Models und PCA.
Diese Modelle haben eine feste Anzahl von Parametern, die aus den Daten gelernt werden.

Nichtparametrische Methoden (nonparametric methods) funktionieren anders.
Sie legen nicht im Voraus eine feste endliche Modellform mit fester Parameterzahl fest.
Stattdessen kann ihre effektive Komplexität mit der Datenmenge wachsen.

\begin{conceptbox}
Nichtparametrische Methoden (nonparametric methods) sind Methoden, deren effektive Modellkomplexität nicht durch eine feste endliche Parameterzahl begrenzt ist.

Sie verwenden oft die Trainingsdaten selbst direkt für Vorhersagen oder Dichteschätzungen.
\end{conceptbox}

Nach diesem Kapitel solltest du folgende Fragen beantworten können:

\begin{itemize}
    \item Was bedeutet nichtparametrisch (nonparametric)?
    \item Was ist der Unterschied zwischen parametrischen und nichtparametrischen Methoden?
    \item Was ist ein Histogramm als Dichteschätzer?
    \item Was ist Kernel Density Estimation?
    \item Was ist ein Kernel?
    \item Was ist die Bandbreite (bandwidth)?
    \item Was ist der Bias-Varianz-Trade-off bei Dichteschätzung?
    \item Wie funktioniert \(k\)-nearest neighbors?
    \item Wie verwendet man \(k\)-NN für Klassifikation?
    \item Wie verwendet man \(k\)-NN für Regression?
    \item Warum ist Feature-Skalierung bei nichtparametrischen Methoden wichtig?
    \item Was sind Vorteile und Nachteile nichtparametrischer Methoden?
    \item Warum leiden nichtparametrische Methoden unter dem Fluch der Dimensionalität?
\end{itemize}

\section{Parametrische und nichtparametrische Modelle}

Ein parametrisches Modell (parametric model) besitzt eine feste Anzahl von Parametern.
Diese Anzahl hängt nicht von der Anzahl der Trainingsdaten ab.

Beispiele:

\begin{itemize}
    \item Lineare Regression mit \(D\) Merkmalen:
    \[
        y(\vec{x},\vec{w})
        =
        \vec{w}^{\top}\vec{x}.
    \]

    \item Logistische Regression:
    \[
        p(t=1\mid\vec{x},\vec{w})
        =
        \sigma(\vec{w}^{\top}\vec{x}).
    \]

    \item Gauß-Verteilung:
    \[
        \mathcal{N}(\vec{x}\mid\vec{\mu},\Sigma).
    \]
\end{itemize}

Bei solchen Modellen wird eine feste Form angenommen.
Die Daten bestimmen nur die Parameter.

\begin{definitionbox}
Ein parametrisches Modell (parametric model) hat eine feste endliche Anzahl von Parametern.

Die Anzahl der Parameter bleibt gleich, auch wenn die Anzahl der Trainingsdaten wächst.
\end{definitionbox}

Nichtparametrische Methoden haben keine feste endliche Parameterzahl in diesem Sinn.
Das bedeutet nicht, dass sie keine Parameter oder Hyperparameter haben.
Sie können Hyperparameter wie Bandbreite oder Anzahl der Nachbarn besitzen.
Aber ihre effektive Komplexität kann mit der Datenmenge wachsen.

\begin{definitionbox}
Eine nichtparametrische Methode (nonparametric method) ist nicht auf eine feste endliche Parameterzahl beschränkt.

Die Modellkomplexität kann mit der Anzahl der Datenpunkte wachsen.
\end{definitionbox}

\begin{notebox}
Nichtparametrisch bedeutet nicht parameterfrei.

Zum Beispiel hat Kernel Density Estimation eine Bandbreite \(h\).
\(k\)-nearest neighbors hat den Hyperparameter \(k\).
Diese Größen sind aber keine festen Modellparameter wie die Gewichte eines linearen Modells.
\end{notebox}

\section{Warum nichtparametrische Methoden?}

Parametrische Modelle sind oft effizient und gut interpretierbar.
Sie können aber zu eingeschränkt sein, wenn die wahre Datenstruktur komplex ist.

Nichtparametrische Methoden sind flexibler.
Sie können komplexe Strukturen abbilden, weil sie weniger starke Annahmen über die Form der Datenverteilung oder Entscheidungsfunktion machen.

\begin{conceptbox}
Nichtparametrische Methoden sind nützlich, wenn man wenig über die Form der zugrunde liegenden Funktion oder Verteilung voraussetzen möchte.
\end{conceptbox}

Typische Anwendungen:

\begin{itemize}
    \item Dichteschätzung (density estimation),
    \item Klassifikation (classification),
    \item Regression (regression),
    \item Glättung von Daten,
    \item explorative Datenanalyse,
    \item einfache Baseline-Modelle.
\end{itemize}

\section{Dichteschätzung}

Bei der Dichteschätzung (density estimation) wollen wir aus Daten eine Wahrscheinlichkeitsdichte schätzen.

Gegeben seien Datenpunkte

\[
    x_1,\dots,x_N
\]

aus einer unbekannten Verteilung mit Dichte

\[
    p(x).
\]

Wir wollen eine Schätzung

\[
    \hat{p}(x)
\]

konstruieren.

\begin{definitionbox}
Dichteschätzung (density estimation) bedeutet, aus Daten eine unbekannte Wahrscheinlichkeitsdichte \(p(x)\) zu schätzen.
\end{definitionbox}

Parametrisch könnten wir zum Beispiel annehmen:

\[
    p(x)
    =
    \mathcal{N}(x\mid\mu,\sigma^2).
\]

Dann müssten wir nur \(\mu\) und \(\sigma^2\) schätzen.
Nichtparametrisch versuchen wir dagegen, die Dichteform flexibler aus den Daten selbst zu rekonstruieren.

\section{Histogramm-Dichteschätzung}

Das Histogramm ist eine einfache nichtparametrische Dichteschätzung.

Wir teilen den Wertebereich in Intervalle, sogenannte Bins.
Für jeden Bin zählen wir, wie viele Datenpunkte hineinfallen.

Sei die Bin-Breite

\[
    \Delta.
\]

Sei \(n_j\) die Anzahl der Datenpunkte im \(j\)-ten Bin.
Dann ist die geschätzte Dichte im \(j\)-ten Bin:

\[
    \hat{p}(x)
    =
    \frac{n_j}{N\Delta}
\]

für \(x\) im \(j\)-ten Bin.

\begin{definitionbox}
Ein Histogramm schätzt die Dichte durch

\[
    \hat{p}(x)
    =
    \frac{\text{Anzahl der Punkte im Bin}}{N\cdot \text{Bin-Breite}}.
\]
\end{definitionbox}

Warum teilen wir durch \(N\Delta\)?

Die Wahrscheinlichkeit für den Bin ist ungefähr:

\[
    \frac{n_j}{N}.
\]

Für eine Dichte gilt näherungsweise:

\[
    p(x)\Delta
    \approx
    \frac{n_j}{N}.
\]

Also:

\[
    p(x)
    \approx
    \frac{n_j}{N\Delta}.
\]

\begin{notebox}
Ein Histogramm ist einfach und anschaulich, aber die Schätzung hängt stark von der Bin-Breite und von der Lage der Bin-Grenzen ab.
\end{notebox}

\section{Bin-Breite und Bias-Varianz-Trade-off}

Die Bin-Breite \(\Delta\) steuert die Glättung.

\subsection{Große Bin-Breite}

Wenn \(\Delta\) groß ist, werden viele Datenpunkte in einem Bin zusammengefasst.
Die Schätzung ist glatt, aber grob.

\[
    \Delta \text{ groß}
    \quad
    \Rightarrow
    \quad
    \text{hoher Bias, niedrige Varianz}.
\]

\subsection{Kleine Bin-Breite}

Wenn \(\Delta\) klein ist, gibt es viele Bins.
Dann schwankt das Histogramm stark, weil manche Bins wenige oder keine Punkte enthalten.

\[
    \Delta \text{ klein}
    \quad
    \Rightarrow
    \quad
    \text{niedriger Bias, hohe Varianz}.
\]

\begin{conceptbox}
Die Bin-Breite eines Histogramms steuert den Bias-Varianz-Trade-off.

Zu große Bins glätten zu stark.
Zu kleine Bins reagieren zu stark auf Zufallsschwankungen.
\end{conceptbox}

\section{Kernel Density Estimation}

Kernel Density Estimation, kurz KDE, ist eine glattere Form der nichtparametrischen Dichteschätzung.

Die Idee:
Statt harte Bins zu zählen, legt man um jeden Datenpunkt eine kleine glatte Dichte.
Die Gesamtschätzung ist der Durchschnitt dieser Dichten.

Für eindimensionale Daten lautet die Schätzung:

\[
    \hat{p}(x)
    =
    \frac{1}{N}
    \sum_{n=1}^{N}
    \frac{1}{h}
    K
    \left(
        \frac{x-x_n}{h}
    \right).
\]

Dabei ist:

\begin{itemize}
    \item \(K\) der Kernel,
    \item \(h>0\) die Bandbreite (bandwidth),
    \item \(x_n\) ein Trainingsdatenpunkt.
\end{itemize}

\begin{definitionbox}
Kernel Density Estimation schätzt eine Dichte durch

\[
    \hat{p}(x)
    =
    \frac{1}{N}
    \sum_{n=1}^{N}
    \frac{1}{h}
    K
    \left(
        \frac{x-x_n}{h}
    \right).
\]
\end{definitionbox}

\begin{conceptbox}
KDE ersetzt jeden Datenpunkt durch eine kleine glatte Glocke.

Die geschätzte Dichte ist die Summe dieser Glocken.
\end{conceptbox}

\section{Kernel}

Ein Kernel \(K(u)\) ist eine Funktion, die typischerweise folgende Eigenschaften hat:

\[
    K(u)\geq 0
\]

und

\[
    \int K(u)\,du=1.
\]

Dadurch ist \(K\) selbst eine Dichte.

Ein häufiger Kernel ist der Gauß-Kernel:

\[
    K(u)
    =
    \frac{1}{\sqrt{2\pi}}
    \exp
    \left(
        -\frac{u^2}{2}
    \right).
\]

Dann ist die KDE:

\[
    \hat{p}(x)
    =
    \frac{1}{N}
    \sum_{n=1}^{N}
    \frac{1}{h\sqrt{2\pi}}
    \exp
    \left(
        -\frac{(x-x_n)^2}{2h^2}
    \right).
\]

\begin{definitionbox}
Der Gauß-Kernel ist:

\[
    K(u)
    =
    \frac{1}{\sqrt{2\pi}}
    \exp
    \left(
        -\frac{u^2}{2}
    \right).
\]
\end{definitionbox}

Andere mögliche Kernel sind:

\begin{itemize}
    \item Rechteck-Kernel,
    \item Dreieck-Kernel,
    \item Epanechnikov-Kernel,
    \item Gauß-Kernel.
\end{itemize}

\begin{notebox}
In der Praxis ist die Wahl der Bandbreite \(h\) meist wichtiger als die genaue Wahl des Kernels.
\end{notebox}

\section{Bandbreite}

Die Bandbreite (bandwidth) \(h\) steuert die Breite des Kernels.

\subsection{Kleine Bandbreite}

Bei kleiner Bandbreite ist jeder Kernel sehr schmal.
Die Dichteschätzung kann stark schwanken.

\[
    h \text{ klein}
    \quad
    \Rightarrow
    \quad
    \text{hohe Varianz}.
\]

\subsection{Große Bandbreite}

Bei großer Bandbreite ist jeder Kernel breit.
Die Dichteschätzung wird sehr glatt.

\[
    h \text{ groß}
    \quad
    \Rightarrow
    \quad
    \text{hoher Bias}.
\]

\begin{conceptbox}
Die Bandbreite \(h\) ist der zentrale Hyperparameter bei Kernel Density Estimation.

Sie kontrolliert den Bias-Varianz-Trade-off.
\end{conceptbox}

\section{Mehrdimensionale KDE}

Für Datenpunkte

\[
    \vec{x}_n\in\mathbb{R}^D
\]

kann KDE verallgemeinert werden.

Mit isotroper Bandbreite \(h\):

\[
    \hat{p}(\vec{x})
    =
    \frac{1}{N}
    \sum_{n=1}^{N}
    \frac{1}{h^D}
    K
    \left(
        \frac{\vec{x}-\vec{x}_n}{h}
    \right).
\]

Für einen multivariaten Gauß-Kernel:

\[
    K(\vec{u})
    =
    \frac{1}{(2\pi)^{D/2}}
    \exp
    \left(
        -\frac{1}{2}
        \vec{u}^{\top}\vec{u}
    \right).
\]

Dann:

\[
    \hat{p}(\vec{x})
    =
    \frac{1}{N}
    \sum_{n=1}^{N}
    \frac{1}{(2\pi h^2)^{D/2}}
    \exp
    \left(
        -\frac{\|\vec{x}-\vec{x}_n\|^2}{2h^2}
    \right).
\]

\begin{definitionbox}
Mehrdimensionale KDE mit Gauß-Kernel ist:

\[
    \hat{p}(\vec{x})
    =
    \frac{1}{N}
    \sum_{n=1}^{N}
    \mathcal{N}(\vec{x}\mid\vec{x}_n,h^2I).
\]
\end{definitionbox}

\section{KDE und Gaussian Mixture Models}

Mehrdimensionale KDE mit Gauß-Kernel kann als spezielles Gaussian Mixture Model interpretiert werden.

\[
    \hat{p}(\vec{x})
    =
    \frac{1}{N}
    \sum_{n=1}^{N}
    \mathcal{N}(\vec{x}\mid\vec{x}_n,h^2I).
\]

Das ist eine Mischung aus \(N\) Gauß-Komponenten.
Jede Komponente hat:

\[
    \vec{\mu}_n=\vec{x}_n,
\]

\[
    \Sigma_n=h^2I,
\]

und Mischgewicht:

\[
    \pi_n=\frac{1}{N}.
\]

\begin{conceptbox}
KDE mit Gauß-Kernel ist wie ein GMM mit einer Komponente pro Datenpunkt.

Die Zentren sind nicht gelernt, sondern direkt die Trainingsdaten.
\end{conceptbox}

\section{\(k\)-nearest neighbors}

Eine weitere zentrale nichtparametrische Methode ist \(k\)-nearest neighbors, kurz \(k\)-NN.

Die Idee:
Für eine neue Eingabe \(\vec{x}\) suchen wir die \(k\) nächsten Trainingspunkte.
Die Vorhersage wird aus deren Zielwerten gebildet.

\begin{definitionbox}
\(k\)-nearest neighbors (\(k\)-NN) sagt für einen neuen Punkt \(\vec{x}\) basierend auf den \(k\) nächsten Trainingspunkten vorher.
\end{definitionbox}

Gegeben sei ein Trainingsdatensatz:

\[
    \mathcal{D}
    =
    \{(\vec{x}_1,t_1),\dots,(\vec{x}_N,t_N)\}.
\]

Für einen neuen Punkt \(\vec{x}\) bestimmen wir die Menge der \(k\) nächsten Nachbarn:

\[
    \mathcal{N}_k(\vec{x})
    =
    \text{Menge der } k \text{ Trainingspunkte mit kleinster Distanz zu } \vec{x}.
\]

\begin{notebox}
\(k\)-NN ist ein lazy learning Verfahren.

Das bedeutet:
Beim Training wird fast nichts berechnet.
Die eigentliche Arbeit passiert bei der Vorhersage.
\end{notebox}

\section{\(k\)-NN für Klassifikation}

Bei Klassifikation sind die Zielwerte Klassen:

\[
    t_n\in\{1,\dots,C\}.
\]

Für einen neuen Punkt \(\vec{x}\) suchen wir die \(k\) nächsten Nachbarn.
Dann wählen wir die Klasse, die unter diesen Nachbarn am häufigsten vorkommt.

\[
    \hat{t}
    =
    \arg\max_{c}
    \sum_{n\in\mathcal{N}_k(\vec{x})}
    \mathbb{1}(t_n=c).
\]

\begin{definitionbox}
\(k\)-NN-Klassifikation verwendet Mehrheitsentscheidung:

\[
    \hat{t}
    =
    \arg\max_{c}
    \sum_{n\in\mathcal{N}_k(\vec{x})}
    \mathbb{1}(t_n=c).
\]
\end{definitionbox}

\begin{examplebox}
Wenn \(k=5\) und die Nachbarn die Klassen

\[
    1,1,2,1,3
\]

haben, dann kommt Klasse \(1\) dreimal vor.
Also ist:

\[
    \hat{t}=1.
\]
\end{examplebox}

\section{\(k\)-NN für Regression}

Bei Regression sind die Zielwerte reell:

\[
    t_n\in\mathbb{R}.
\]

Für einen neuen Punkt \(\vec{x}\) wird meist der Durchschnitt der Zielwerte der \(k\) nächsten Nachbarn verwendet:

\[
    \hat{y}(\vec{x})
    =
    \frac{1}{k}
    \sum_{n\in\mathcal{N}_k(\vec{x})}
    t_n.
\]

\begin{definitionbox}
\(k\)-NN-Regression verwendet den Mittelwert der Zielwerte der \(k\) nächsten Nachbarn:

\[
    \hat{y}(\vec{x})
    =
    \frac{1}{k}
    \sum_{n\in\mathcal{N}_k(\vec{x})}
    t_n.
\]
\end{definitionbox}

Man kann auch gewichtete \(k\)-NN-Regression verwenden.
Dann bekommen nähere Punkte ein größeres Gewicht:

\[
    \hat{y}(\vec{x})
    =
    \frac{
        \sum_{n\in\mathcal{N}_k(\vec{x})}
        w_n t_n
    }{
        \sum_{n\in\mathcal{N}_k(\vec{x})}
        w_n
    }.
\]

Ein mögliches Gewicht ist:

\[
    w_n
    =
    \frac{1}{\|\vec{x}-\vec{x}_n\|+\varepsilon}.
\]

\section{Wahl von \(k\)}

Der Hyperparameter \(k\) steuert die Glättung.

\subsection{Kleines \(k\)}

Wenn \(k=1\), wird der Zielwert des nächsten Trainingspunkts verwendet.
Das Modell kann sehr flexibel sein, aber stark auf Rauschen reagieren.

\[
    k \text{ klein}
    \quad
    \Rightarrow
    \quad
    \text{niedriger Bias, hohe Varianz}.
\]

\subsection{Großes \(k\)}

Wenn \(k\) groß ist, werden viele Nachbarn gemittelt.
Das Modell ist glatter, kann aber lokale Struktur verlieren.

\[
    k \text{ groß}
    \quad
    \Rightarrow
    \quad
    \text{hoher Bias, niedrige Varianz}.
\]

\begin{conceptbox}
Bei \(k\)-NN steuert \(k\) den Bias-Varianz-Trade-off.

Kleines \(k\): flexibel, aber rauschsensitiv.

Großes \(k\): stabil, aber möglicherweise zu glatt.
\end{conceptbox}

\section{Entscheidungsgrenzen bei \(k\)-NN}

Bei \(1\)-NN können Entscheidungsgrenzen sehr komplex sein.
Jeder Trainingspunkt beeinflusst eine Region des Raums.
Diese Regionen hängen mit Voronoi-Zellen zusammen.

\begin{definitionbox}
Eine Voronoi-Zelle eines Trainingspunkts enthält alle Punkte im Raum, die näher an diesem Trainingspunkt liegen als an allen anderen Trainingspunkten.
\end{definitionbox}

Bei \(1\)-NN wird jeder Punkt im Raum der Klasse des nächstgelegenen Trainingspunkts zugeordnet.
Dadurch können sehr unregelmäßige Entscheidungsgrenzen entstehen.

Wenn \(k\) größer wird, werden die Entscheidungsgrenzen glatter.

\begin{notebox}
\(1\)-NN hat oft sehr niedrigen Trainingsfehler, kann aber stark überanpassen.
\end{notebox}

\section{Distanzmaße}

Die Wahl des Distanzmaßes ist entscheidend.

\subsection{Euklidische Distanz}

\[
    d(\vec{x},\vec{x}')
    =
    \sqrt{
        \sum_{j=1}^{D}
        (x_j-x_j')^2
    }.
\]

\subsection{Manhattan-Distanz}

\[
    d(\vec{x},\vec{x}')
    =
    \sum_{j=1}^{D}
    |x_j-x_j'|.
\]

\subsection{Minkowski-Distanz}

\[
    d_p(\vec{x},\vec{x}')
    =
    \left(
        \sum_{j=1}^{D}
        |x_j-x_j'|^p
    \right)^{1/p}.
\]

Für \(p=2\) erhält man die euklidische Distanz.
Für \(p=1\) erhält man die Manhattan-Distanz.

\begin{definitionbox}
Die Minkowski-Distanz ist:

\[
    d_p(\vec{x},\vec{x}')
    =
    \left(
        \sum_{j=1}^{D}
        |x_j-x_j'|^p
    \right)^{1/p}.
\]
\end{definitionbox}

\begin{notebox}
Das richtige Distanzmaß hängt vom Datentyp und Problem ab.
Für Bilder, Texte oder kategoriale Daten sind einfache euklidische Distanzen nicht immer sinnvoll.
\end{notebox}

\section{Feature-Skalierung bei \(k\)-NN}

Da \(k\)-NN direkt auf Distanzen basiert, ist Feature-Skalierung sehr wichtig.

Wenn ein Merkmal eine viel größere Skala hat als andere, dominiert es die Distanz.

Eine Standardisierung ist:

\[
    x_j'
    =
    \frac{x_j-\mu_j}{\sigma_j}.
\]

Dabei sind \(\mu_j\) und \(\sigma_j\) Mittelwert und Standardabweichung des Merkmals im Trainingsdatensatz.

\begin{examtipbox}
Bei \(k\)-NN und KDE sollte man Merkmale mit unterschiedlichen Einheiten oder Skalen fast immer skalieren oder standardisieren.

Die Skalierungsparameter müssen nur auf den Trainingsdaten berechnet werden, um Data Leakage zu vermeiden.
\end{examtipbox}

\section{Fluch der Dimensionalität}

Nichtparametrische Methoden leiden stark unter dem Fluch der Dimensionalität (curse of dimensionality).

In hohen Dimensionen werden Datenpunkte sehr dünn verteilt.
Lokale Nachbarschaften enthalten dann oft kaum noch sinnvolle Information.

Bei \(k\)-NN bedeutet das:
Der nächste Nachbar kann trotzdem weit entfernt sein.
Dann ist die Annahme lokaler Ähnlichkeit schwächer.

Bei KDE bedeutet das:
Man braucht sehr viele Datenpunkte, um eine Dichte in hoher Dimension zuverlässig zu schätzen.

\begin{conceptbox}
Nichtparametrische Methoden brauchen in hoher Dimension sehr viele Daten.

Der Grund ist:
Lokale Methoden benötigen genügend Datenpunkte in der Nähe jedes Abfragepunkts.
\end{conceptbox}

\section{Speicher- und Rechenaufwand}

Viele nichtparametrische Methoden speichern die Trainingsdaten.

Bei \(k\)-NN muss man für eine neue Eingabe oft Distanzen zu vielen Trainingspunkten berechnen.

Naiv kostet eine Vorhersage:

\[
    O(ND),
\]

wobei:

\begin{itemize}
    \item \(N\) die Anzahl der Trainingsdaten,
    \item \(D\) die Dimension ist.
\end{itemize}

\begin{notebox}
\(k\)-NN ist beim Training billig, aber bei der Vorhersage teuer.

Parametrische Modelle sind oft beim Training teurer, aber bei der Vorhersage schneller.
\end{notebox}

Zur Beschleunigung kann man Datenstrukturen verwenden:

\begin{itemize}
    \item \(k\)-d-Trees,
    \item Ball Trees,
    \item approximate nearest neighbor search.
\end{itemize}

In sehr hoher Dimension funktionieren klassische Baumstrukturen aber oft schlechter.

\section{Parzen-Fenster}

Kernel Density Estimation wird auch als Parzen-Fenster-Methode (Parzen window method) bezeichnet.

Die Grundidee:
Ein Fenster wird um den Punkt \(x\) gelegt.
Man zählt oder gewichtet, wie viele Trainingspunkte in diesem Fenster liegen.

Mit einem Rechteck-Kernel zählt man Punkte innerhalb eines festen Fensters.
Mit einem Gauß-Kernel werden alle Punkte berücksichtigt, aber nahe Punkte stärker gewichtet.

\begin{definitionbox}
Parzen-Fenster-Dichteschätzung ist eine nichtparametrische Dichteschätzung, bei der lokale Nachbarschaften zur Schätzung von \(p(x)\) verwendet werden.
\end{definitionbox}

\section{\(k\)-NN-Dichteschätzung}

Neben KDE gibt es auch \(k\)-NN-Dichteschätzung.
Die Idee:
Für einen Punkt \(\vec{x}\) sucht man das kleinste Volumen \(V\), das \(k\) Trainingspunkte enthält.

Wenn \(N\) Datenpunkte insgesamt vorhanden sind, ist die geschätzte Wahrscheinlichkeit in diesem Volumen:

\[
    \frac{k}{N}.
\]

Die Dichte ist ungefähr Wahrscheinlichkeit pro Volumen:

\[
    \hat{p}(\vec{x})
    =
    \frac{k}{NV}.
\]

\begin{definitionbox}
Die \(k\)-NN-Dichteschätzung ist:

\[
    \hat{p}(\vec{x})
    =
    \frac{k}{N V(\vec{x})},
\]

wobei \(V(\vec{x})\) das Volumen der kleinsten Umgebung um \(\vec{x}\) ist, die \(k\) Datenpunkte enthält.
\end{definitionbox}

\begin{conceptbox}
KDE verwendet feste Bandbreite und variable Punktanzahl.

\(k\)-NN-Dichteschätzung verwendet feste Punktanzahl und variable Bandbreite.
\end{conceptbox}

\section{Nichtparametrische Regression als lokale Mittelung}

\(k\)-NN-Regression kann als lokale Mittelung verstanden werden.

Für einen Punkt \(\vec{x}\) wird nicht ein globaler Parametervektor gelernt.
Stattdessen wird lokal in der Umgebung von \(\vec{x}\) gemittelt:

\[
    \hat{y}(\vec{x})
    =
    \frac{1}{k}
    \sum_{n\in\mathcal{N}_k(\vec{x})}
    t_n.
\]

Allgemeiner kann man Kernel-Regression verwenden:

\[
    \hat{y}(\vec{x})
    =
    \frac{
        \sum_{n=1}^{N}
        K
        \left(
            \frac{\|\vec{x}-\vec{x}_n\|}{h}
        \right)
        t_n
    }{
        \sum_{n=1}^{N}
        K
        \left(
            \frac{\|\vec{x}-\vec{x}_n\|}{h}
        \right)
    }.
\]

\begin{definitionbox}
Kernel-Regression sagt einen Zielwert als gewichteten Durchschnitt der Trainingszielwerte vorher:

\[
    \hat{y}(\vec{x})
    =
    \frac{
        \sum_{n=1}^{N}
        w_n(\vec{x})t_n
    }{
        \sum_{n=1}^{N}
        w_n(\vec{x})
    }.
\]

Die Gewichte \(w_n(\vec{x})\) sind größer für Trainingspunkte nahe bei \(\vec{x}\).
\end{definitionbox}

\section{Vergleich parametrisch und nichtparametrisch}

\[
\begin{array}{c|c|c}
\text{Eigenschaft}
&
\text{parametrisch}
&
\text{nichtparametrisch}
\\
\hline
\text{Parameterzahl}
&
\text{fest}
&
\text{kann mit Daten wachsen}
\\
\hline
\text{Training}
&
\text{oft Optimierung}
&
\text{oft Speichern der Daten}
\\
\hline
\text{Vorhersage}
&
\text{oft schnell}
&
\text{oft teuer}
\\
\hline
\text{Flexibilität}
&
\text{durch Modellform begrenzt}
&
\text{hoch}
\\
\hline
\text{Datenbedarf}
&
\text{oft geringer}
&
\text{oft höher}
\\
\hline
\text{Interpretierbarkeit}
&
\text{modellabhängig}
&
\text{oft lokal}
\end{array}
\]

\begin{conceptbox}
Parametrische Modelle verwenden starke Annahmen und können dadurch effizient lernen.

Nichtparametrische Methoden machen weniger starke Annahmen, brauchen aber oft mehr Daten.
\end{conceptbox}

\section{Vorteile nichtparametrischer Methoden}

\begin{itemize}
    \item flexibel,
    \item einfach zu verstehen,
    \item wenige Modellannahmen,
    \item gute Baselines,
    \item natürlich lokal,
    \item können komplexe Entscheidungsgrenzen modellieren.
\end{itemize}

\begin{conceptbox}
Nichtparametrische Methoden sind besonders nützlich, wenn lokale Ähnlichkeit aussagekräftig ist.
\end{conceptbox}

\section{Nachteile nichtparametrischer Methoden}

\begin{itemize}
    \item hoher Speicherbedarf,
    \item teure Vorhersage bei großen Datensätzen,
    \item empfindlich gegenüber Feature-Skalierung,
    \item empfindlich gegenüber irrelevanten Merkmalen,
    \item leiden stark unter hoher Dimension,
    \item Hyperparameter wie \(k\) oder \(h\) müssen gewählt werden,
    \item weniger kompakte Modellzusammenfassung.
\end{itemize}

\begin{notebox}
Nichtparametrische Methoden sind nicht automatisch besser, nur weil sie flexibler sind.

Ohne genug Daten oder sinnvolle Distanzen können sie schlechter generalisieren als einfache parametrische Modelle.
\end{notebox}

\section{Modellwahl und Validierung}

Hyperparameter nichtparametrischer Methoden sollten mit Validierungsdaten gewählt werden.

Beispiele:

\[
    k
\]

bei \(k\)-NN,

\[
    h
\]

bei KDE oder Kernel-Regression,

\[
    \Delta
\]

beim Histogramm.

Man kann wählen:

\[
    k^*
    =
    \arg\min_k
    E_{\mathrm{val}}(k),
\]

oder

\[
    h^*
    =
    \arg\min_h
    E_{\mathrm{val}}(h).
\]

\begin{examtipbox}
Auch bei nichtparametrischen Methoden gilt:

Testdaten dürfen nicht zur Wahl von \(k\), \(h\), Distanzmaß oder Vorverarbeitung verwendet werden.
\end{examtipbox}

\section{Mini-Aufgaben}

\subsection{Aufgabe 1: Histogramm-Dichte}

Ein Datensatz enthält

\[
    N=100
\]

Punkte.
In einem Bin der Breite

\[
    \Delta=0{,}5
\]

liegen

\[
    n_j=10
\]

Punkte.
Berechne die Histogramm-Dichteschätzung in diesem Bin.

\begin{examplebox}
Die Histogramm-Dichte ist:

\[
    \hat{p}(x)
    =
    \frac{n_j}{N\Delta}.
\]

Einsetzen:

\[
    \hat{p}(x)
    =
    \frac{10}{100\cdot 0{,}5}
    =
    \frac{10}{50}
    =
    0{,}2.
\]
\end{examplebox}

\subsection{Aufgabe 2: \(k\)-NN-Klassifikation}

Für einen neuen Punkt seien die \(5\) nächsten Nachbarn mit Klassen:

\[
    A,A,B,A,B.
\]

Welche Klasse sagt \(5\)-NN vorher?

\begin{examplebox}
Klasse \(A\) kommt dreimal vor.
Klasse \(B\) kommt zweimal vor.

Also sagt \(5\)-NN voraus:

\[
    \hat{t}=A.
\]
\end{examplebox}

\subsection{Aufgabe 3: \(k\)-NN-Regression}

Die \(3\) nächsten Nachbarn haben Zielwerte:

\[
    2,
    \quad
    4,
    \quad
    10.
\]

Berechne die \(3\)-NN-Regressionsvorhersage.

\begin{examplebox}
Die Vorhersage ist der Mittelwert:

\[
    \hat{y}
    =
    \frac{2+4+10}{3}
    =
    \frac{16}{3}.
\]

Also:

\[
    \hat{y}
    \approx
    5{,}33.
\]
\end{examplebox}

\subsection{Aufgabe 4: Bias und Varianz bei \(k\)}

Was passiert typischerweise, wenn \(k\) bei \(k\)-NN sehr klein gewählt wird?

\begin{examplebox}
Ein sehr kleines \(k\), zum Beispiel \(k=1\), macht das Modell sehr flexibel.

Typischerweise gilt:

\[
    \text{Bias niedrig},
    \qquad
    \text{Varianz hoch}.
\]

Das Modell kann stark auf Rauschen und einzelne Trainingspunkte reagieren.
\end{examplebox}

\subsection{Aufgabe 5: KDE als GMM}

Erkläre, warum KDE mit Gauß-Kernel wie ein Gaussian Mixture Model mit \(N\) Komponenten aussieht.

\begin{examplebox}
KDE mit Gauß-Kernel ist:

\[
    \hat{p}(\vec{x})
    =
    \frac{1}{N}
    \sum_{n=1}^{N}
    \mathcal{N}(\vec{x}\mid\vec{x}_n,h^2I).
\]

Das ist eine Mischung von \(N\) Gauß-Verteilungen.

Jede Komponente hat:

\[
    \vec{\mu}_n=\vec{x}_n,
\]

\[
    \Sigma_n=h^2I,
\]

und

\[
    \pi_n=\frac{1}{N}.
\]

Daher kann man KDE als spezielles GMM mit einer Komponente pro Datenpunkt interpretieren.
\end{examplebox}

\section{Häufige Fehler}

\begin{examtipbox}
Typische Fehler bei nichtparametrischen Methoden:

\begin{enumerate}
    \item Nichtparametrisch mit parameterfrei verwechseln.
    \item Histogramme ohne Beachtung der Bin-Breite interpretieren.
    \item Bei KDE die Bandbreite \(h\) ignorieren.
    \item Bei \(k\)-NN die Feature-Skalierung vergessen.
    \item Euklidische Distanz verwenden, obwohl sie für die Daten nicht sinnvoll ist.
    \item \(k=1\) verwenden und Overfitting übersehen.
    \item Sehr großes \(k\) verwenden und Underfitting übersehen.
    \item Testdaten zur Wahl von \(k\) oder \(h\) verwenden.
    \item Den Fluch der Dimensionalität unterschätzen.
    \item \(k\)-NN als schnelles Vorhersagemodell bei riesigen Datensätzen erwarten.
    \item Irrelevante Merkmale nicht entfernen, obwohl sie Distanzen verzerren.
    \item KDE in hoher Dimension ohne genug Daten verwenden.
\end{enumerate}
\end{examtipbox}

\section{Zusammenfassung}

\begin{itemize}
    \item Parametrische Modelle haben eine feste endliche Parameterzahl.
    \item Nichtparametrische Methoden haben keine feste endliche Parameterzahl in diesem Sinn.
    \item Nichtparametrisch bedeutet nicht parameterfrei.
    \item Histogramme sind einfache nichtparametrische Dichteschätzer.
    \item Die Histogramm-Dichte in einem Bin ist:
    \[
        \hat{p}(x)
        =
        \frac{n_j}{N\Delta}.
    \]
    \item Kernel Density Estimation schätzt:
    \[
        \hat{p}(x)
        =
        \frac{1}{N}
        \sum_{n=1}^{N}
        \frac{1}{h}
        K
        \left(
            \frac{x-x_n}{h}
        \right).
    \]
    \item Die Bandbreite \(h\) kontrolliert Glättung, Bias und Varianz.
    \item Mehrdimensionale KDE mit Gauß-Kernel ist:
    \[
        \hat{p}(\vec{x})
        =
        \frac{1}{N}
        \sum_{n=1}^{N}
        \mathcal{N}(\vec{x}\mid\vec{x}_n,h^2I).
    \]
    \item \(k\)-NN verwendet die \(k\) nächsten Trainingspunkte für Vorhersagen.
    \item \(k\)-NN-Klassifikation verwendet Mehrheitsentscheidung:
    \[
        \hat{t}
        =
        \arg\max_c
        \sum_{n\in\mathcal{N}_k(\vec{x})}
        \mathbb{1}(t_n=c).
    \]
    \item \(k\)-NN-Regression verwendet lokale Mittelung:
    \[
        \hat{y}(\vec{x})
        =
        \frac{1}{k}
        \sum_{n\in\mathcal{N}_k(\vec{x})}
        t_n.
    \]
    \item Kleines \(k\) bedeutet niedrigen Bias und hohe Varianz.
    \item Großes \(k\) bedeutet höheren Bias und niedrigere Varianz.
    \item Nichtparametrische Methoden sind flexibel, brauchen aber oft viele Daten.
    \item Feature-Skalierung und Distanzwahl sind entscheidend.
    \item In hoher Dimension leiden nichtparametrische Methoden stark unter dem Fluch der Dimensionalität.
\end{itemize}

\section{Typische Prüfungsfragen}

\begin{examtipbox}
Mögliche Prüfungsfragen zu diesem Kapitel:

\begin{enumerate}
    \item Was bedeutet nichtparametrisch?
    \item Warum heißt nichtparametrisch nicht parameterfrei?
    \item Was ist der Unterschied zwischen parametrischen und nichtparametrischen Modellen?
    \item Was ist Dichteschätzung?
    \item Wie funktioniert ein Histogramm als Dichteschätzer?
    \item Welche Rolle spielt die Bin-Breite?
    \item Was ist Kernel Density Estimation?
    \item Was ist ein Kernel?
    \item Was ist die Bandbreite \(h\)?
    \item Wie beeinflusst \(h\) Bias und Varianz?
    \item Schreibe KDE mit Gauß-Kernel auf.
    \item Wie hängt KDE mit Gaussian Mixture Models zusammen?
    \item Was ist \(k\)-nearest neighbors?
    \item Wie funktioniert \(k\)-NN-Klassifikation?
    \item Wie funktioniert \(k\)-NN-Regression?
    \item Welche Rolle spielt der Hyperparameter \(k\)?
    \item Warum ist \(1\)-NN oft varianzreich?
    \item Warum kann großes \(k\) zu Underfitting führen?
    \item Welche Distanzmaße können bei \(k\)-NN verwendet werden?
    \item Warum ist Feature-Skalierung bei \(k\)-NN wichtig?
    \item Was ist \(k\)-NN-Dichteschätzung?
    \item Was ist Kernel-Regression?
    \item Warum sind nichtparametrische Methoden in hoher Dimension schwierig?
    \item Welche Vor- und Nachteile haben nichtparametrische Methoden?
\end{enumerate}
\end{examtipbox}